\section{aim}

the task is to classify the emotion expressed in a piece of text with a deep learning model. the usual way to
do this is to fine-tune a BERT-style model on a six-class twitter dataset (dair-ai/emotion) and report
accuracy, which ends up around 93\%. i wanted something harder and closer to how emotion actually shows up
in text, so i used GoEmotions \cite{goemotions}: 27 fine-grained emotions plus neutral, and a comment can
have several of them at once. that makes it a multi-label problem with very imbalanced labels (grief has 77
training examples, neutral has 14k).

the goals were: (1) build a ladder of models from a linear baseline up to a large transformer, (2) get a
small model that is close to the large one but small enough to run in a web browser, (3) make the
probabilities it outputs actually mean something (calibration), (4) check how well it transfers to a
different domain without any extra training, and (5) build a demo that is useful for poking at the model,
including a per-word explanation and an ``emotion arc'' over longer text.

\section{background}

\textbf{GoEmotions.} \cite{goemotions} collected 58k reddit comments, each labelled by three or more raters
with any subset of 27 emotions or neutral. the ``simplified'' release keeps labels that at least two raters
agreed on. their BERT-base baseline gets 0.46 macro-f1 on the test set, and they also group the labels into
Ekman's six basic emotions \cite{ekman} plus neutral, where BERT gets 0.64.

\textbf{multi-label classification.} each label gets its own sigmoid output and binary cross-entropy loss.
the decision threshold does not have to be 0.5: for rare labels the model rarely becomes confident, so a
lower threshold picked on the validation set usually gives a much better f1. asymmetric loss (ASL)
\cite{asl} is a variant of focal loss for multi-label problems that down-weights the many easy negatives
more than the positives.

\textbf{knowledge distillation} \cite{hinton} trains a small student to match the softened outputs of a big
teacher as well as the gold labels. for sigmoid outputs i used the binary version: binary cross-entropy
between $\sigma(z_s/T)$ and $\sigma(z_t/T)$, scaled by $T^2$.

\textbf{calibration.} a classifier is calibrated if, among all the times it says 70\%, the label is
actually present about 70\% of the time. expected calibration error (ECE) measures the gap using binned
predictions. temperature scaling \cite{guo} divides all logits by a single scalar fitted on the validation
set; it does not change any ranking, so f1 at tuned thresholds is unaffected.

\section{data}

\begin{table}[H]
\centering\small
\begin{tabular}{lrrl}
\toprule
dataset & train & val & test \\
\midrule
GoEmotions simplified (reddit, 28 labels, multi-label) & 43\,410 & 5\,426 & 5\,427 \\
dair-ai/emotion (twitter, 6 classes, single label), ood test only & -- & -- & 2\,000 \\
\bottomrule
\end{tabular}
\caption{datasets. i only use the official splits.}
\end{table}

GoEmotions comments are short (12.8 words on average). 16\% of training comments have more than one label.
names and religions are masked as \texttt{[NAME]} and \texttt{[RELIGION]}. the label distribution is
extremely skewed: neutral appears in a third of the comments while grief, pride, relief and nervousness have
fewer than 170 training examples each and 6--23 test examples, so their test f1 is noisy.

for the out-of-domain test i map both label sets to Ekman groups. dair-ai/emotion has sadness, joy, love,
anger, fear and surprise; love goes into joy (GoEmotions also puts love under joy). it has no disgust and no
neutral, so for this test the GoEmotions model's Ekman probabilities are restricted to the five shared
groups and the argmax is compared with the gold class. the twitter text is lowercased with punctuation
stripped, which is quite different from reddit, and none of it was seen in training.
